\documentclass[a4paper,
  DIV=14, % Faktor Nutzefläche vs. Rand (zwischen 6 und 15)
  BCOR=8mm, % Bindekorrekturfaktor für Klebebindung
  twoside,
  english,
  titlepage=false,
  footsepline=true,
  toc=bib,indent,
  fontsize=12pt % Standardschriftgröße, default = 11pt
]{scrbook}
\usepackage[utf8]{inputenc}


%  metadata in Settings.tex



% load settings
%!TEX root =  ./Paper.tex
%!TEX encoding = UTF-8 Unicode

\def\thisversion{4.6}
\def\thisdate{\today}
\def\thistime{\currenttime}

\def\thistitle{IXSI - Interface for X-Sharing Information\\ Version 4}
\def\thistitleshort{IXSI}
\def\thisauthor{
  Markus Beutel,
  Sevket Gökay,
  Peter von Grumbkow,
  Dirk Hillbrecht,
  Karl-Heinz Krempels,
  Christian Samsel,
  Christoph Terwelp,
  Heike Twele,
  Gerhard Wagner}
\def\thispublishers{Cantamen GmbH, HaCon Ing.-Ges. mbH, RWTH Aachen}

\def\thisabstract{
  IXSI ist unter der Creative-Commons-Lizenz \blockquote{Namensnennung Keine Bearbeitung 3.0 Deutschland}, CC BY-ND 3.0 DE, lizensiert. Eine Nutzung der Schnittstellenspezifikation ist frei möglich, auch in anderen Projekten. Bei einer Veröffentlichung müssen die Autoren genannt werden. Eine Abänderung der Spezifikation ist nicht gestattet. Die Details zur Lizenz finden sich in Kurzform (\url{http://creativecommons.org/licenses/by-nd/3.0/de/}) und in der rechtlich bindenden Vertragsform (\url{http://creativecommons.org/licenses/by-nd/3.0/de/legalcode}) auf den Seiten der Creative Commons. Das zugehörige XML Schema kann von \url{http://www.ixsi-schnittstelle.de} herruntergeladen werden.
}
%    
%----------------------------------------------------------
%---- Packages --------------------------------------------
%----------------------------------------------------------



%---- Have to be loaded at first---------------------------
\usepackage{commandsrwth}
\usepackage{etex}
\usepackage{ifthen}
\usepackage{ifpdf}
\usepackage{xspace}
\usepackage{ragged2e}
\usepackage{epsfig}
% \usepackage{subfigure}

\usepackage{calc}
\usepackage{pslatex}

%---- Coding, Language, Fonts, and Symbols ----------------
\usepackage{babel}
\usepackage[utf8]{inputenc}
\usepackage[T1]{fontenc}
\usepackage{textcomp}
\usepackage{eurosym}
\usepackage[expansion=false]{microtype}

%---- Mathematic ------------------------------------------
\usepackage{amsmath}
\usepackage{amssymb}
\usepackage{amsfonts}
\usepackage{amstext}
\usepackage{amsthm}
\usepackage{thmtools}
\usepackage[all]{onlyamsmath}
\usepackage{latexsym}
\usepackage{stmaryrd}
\usepackage{mathtools}
\usepackage{fixmath}
\usepackage{icomma}
\usepackage{cleveref}
%---- Scientific ------------------------------------------
\usepackage{listings}
\usepackage{bytefield}
\usepackage[
separate-uncertainty,
multi-part-units=repeat,
product-units=repeat,
per-mode=symbol-or-fraction,
]
{siunitx}

%---- Text and Layout -------------------------------------
\usepackage[
babel,
threshold=2,
strict
]{csquotes}
\usepackage{upquote}
\usepackage{varwidth}
\usepackage{multicol}


% \usepackage{enumitem}
% within paragraphs
\usepackage[svgnames]{xcolor}


%---- Floats - Tables, Figures ----------------------------

\usepackage{graphicx}
\usepackage{array}
\usepackage{tabularx}
\usepackage{booktabs}
\usepackage{multirow}
\usepackage{floatrow}
\floatsetup[widefigure]{}
\usepackage{longtable}
\usepackage{rotating}
\usepackage{dcolumn}
\usepackage{float}
\floatstyle{plaintop}
\restylefloat{table}
% flafter forces float to be placed after referenced
% \usepackage{flafter} 

%  fancy captions (use default for insticc)
% \usepackage[small,hang]{caption}

%table caption on top
\usepackage[small]{caption}

\usepackage[
caption=false,
]{subfig}

\graphicspath{{./graphics/}}


% permissive float placement
\IfNotClass{beamer}{%
  \setcounter{topnumber}{4}
  \setcounter{bottomnumber}{4}
  \setcounter{totalnumber}{10}
  \renewcommand{\textfraction}{0.15}
  \renewcommand{\topfraction}{0.85}
  \renewcommand{\bottomfraction}{0.70}
  \renewcommand{\floatpagefraction}{0.66}
}


%---- misc ----------------------------------------------

\IfNotClass{beamer}{
  \usepackage[noadjust]{cite}
  \usepackage{paralist}
  \usepackage[hyphens]{url}
  \usepackage{imakeidx}                    % Indexerstellung 
  \usepackage{ftnright} % Fussnoten nur rechts
}

\IfClass{beamer}{%
  \usepackage[accumulated]{beamerseminar}
  \usepackage[absolute,overlay]{textpos}
  \usepackage{pgfpages}                   % Ermöglicht eine Two-Screen Präsentation
  \usepackage{textpos}                    % Absolute Positionierung von Texten/Graphiken
}

% balance columns use for IEEE/ACM
% \usepackage{flushend} 
% \usepackage{balance} 


%---- hyperref --------------------------------------------
\IfNotClass{beamer}{%
  \usepackage[
  %   draft,
  %   final,
  breaklinks=true,
  linktocpage=true,
  hyperfootnotes=true,
%   pagebackref=true,
  plainpages=false,
  naturalnames=true,
  pageanchor=true,
  bookmarks=true,
%   hidelinks,
  colorlinks=false,
  linkcolor=blue,
  urlcolor=blue,
  citecolor=blue,
  bookmarks=true,
  bookmarksnumbered=true,
  pdfpagemode=UseOutlines,
  pdfpagelabels=true,
  pdfencoding=auto,
%  psdextra,
  pdfstartview={FitH},
  unicode
  ]{hyperref}

% not usefull for pdflatex
%\usepackage[hyphenbreaks]{breakurl}

%---- glossaries ------------------------------------------
\usepackage[
  acronym,
  shortcuts,
  style=super,
  description,
  nonumberlist
]{glossaries}
% \renewcommand*{\glspostdescription}{}


} %if Until here

%----------------------------------------------------------
%---- Settings --------------------------------------------
%----------------------------------------------------------


    
%--- hyperref meta data -----------------------------------
\hypersetup{%
%   pdfauthor={\author},
% %   pdfkeywords={\thiskeywords},
%   pdftitle={\title},
}





%---- listings --------------------------------------------
\lstset{%
  language=Java,
  tabsize=2,
  basewidth=0.55em,
  frame=lines,
  framerule=\heavyrulewidth,
  basicstyle=\scriptsize\ttfamily,
  keywordstyle=\bfseries\color{Maroon},
  stringstyle=\color{Maroon},
  commentstyle=\itshape\color{Green},
  showstringspaces=false,
  aboveskip=8pt,
  belowskip=8pt,
  numbers=left,
  numberstyle=\tiny,
  stepnumber=1,
  numbersep=0.5em,
  % xleftmargin=0.5em,
  % xrightmargin=0.5em,
  morecomment=[l]{//},
  escapeinside={(*}{*)},
  captionpos=b,
}

\lstdefinestyle{XML-style} {
  language=XML,
  morestring=[b]",
  %morestring=[s]{>}{<},
  morecomment=[s]{<?}{?>},
  morecomment=[s]{<!--}{-->},
  stringstyle=\color{black},
  tagstyle=\color{Purple},
  keywordstyle=\color{MediumBlue},
  commentstyle=\color{gray}\slshape,
  morekeywords={xmlns:xsi,xmlns:xsd,xmlns:soap,xmlns,version,type,name,minOccurs,maxOccurs,value,base,
    elementFormDefault,targetNamespace,sex}% list your attributes here
  }
  
  \lstdefinestyle{JSON-style} {
    stringstyle=\color{black},
  }
  
  \lstdefinestyle{XMLX-style} {
    backgroundcolor=\color{white},
    frame=leftline,
    language=XML,
    morestring=[b]",
    %morestring=[s]{>}{<},
    morecomment=[s]{<?}{?>},
    morecomment=[s]{<!--}{-->},
    stringstyle=\color{black},
    tagstyle=\color{Purple},
    keywordstyle=\color{MediumBlue},
    commentstyle=\color{gray}\slshape,
    morekeywords={xmlns:xsi,xmlns:xsd,xmlns:soap,xmlns,version,type,name,minOccurs,maxOccurs,value,base,
      elementFormDefault,targetNamespace,sex},% list your attributes here
      numbers=none
    }
    
    \lstdefinestyle{JSONX-style} {
      backgroundcolor=\color{white},
      frame=leftline,
      stringstyle=\color{black},
      numbers=none
    }


%----------------------------------------------------------
%---- Definitions -----------------------------------------
%----------------------------------------------------------

%---- General definitions ---------------------------------

% Margin notes
\providecommand{\marginline}[1]{\marginpar[\RaggedLeft{#1}]{\RaggedRight #1}}
\newcommand*{\TODO}[1]{\marginline{\scriptsize\textcolor{red}{TODO: #1}}}

% Style abbreviations
\newcommand*{\ebf}[1]{\emph{\textbf{#1}}}
\newcommand*{\bs}[1]{\boldsymbol{#1}}
\renewcommand*{\_}{\ensuremath{\mathunderscore}\xspace}

% Enumeration
\newcommand*{\st}{\ensuremath{1^{\mbox{\scriptsize st}}}\xspace}
\newcommand*{\nd}{\ensuremath{2^{\mbox{\scriptsize nd}}}\xspace}
\newcommand*{\rd}{\ensuremath{3^{\mbox{\scriptsize rd}}}\xspace}
\newcommand*{\nth}[1]{\ensuremath{#1^{\mbox{\scriptsize th}}}\xspace}

% Arrows
\renewcommand*{\implies}{\ensuremath{\rightarrow}\xspace}
\renewcommand*{\iff}{\ensuremath{\leftrightarrow}\xspace}
\newcommand*{\IF}{\ensuremath{\Rightarrow}\xspace}
\newcommand*{\sra}{\ensuremath{\shortrightarrow}\xspace}
\newcommand*{\sla}{\ensuremath{\shortleftarrow}\xspace}


\newcolumntype{M}[1]{
  >{\RaggedRight\arraybackslash\hspace{0pt}}m{#1}%
}
\newcolumntype{R}[1]{
  >{\begin{turn}{90}\begin{minipage}{#1}%
  \RaggedRight\arraybackslash\hspace{0pt}}l
  <{\end{minipage}\end{turn}}%
}
\newcommand*{\armultirow}[3]{%
  \multicolumn{#1}{#2}{%
    \begin{picture}(0,0)%
    \put(0,0){\begin{tabular}[t]{@{}#2@{}}#3\end{tabular}}%
    \end{picture}%
  }%
}
\makeatletter
\newcolumntype{K}[1]{%
  >{\DC@{.}{.}{#1}}l<{\DC@end}%
}
\makeatother

%---- hyperref --------------------------------------------
% Write Section instead of section for \ref
\addto\extrasenglish{%
  \renewcommand*{\sectionautorefname}{Section}%
  \renewcommand*{\subsectionautorefname}{\sectionautorefname}%
  \renewcommand*{\subsubsectionautorefname}{\sectionautorefname}%
}

\newcommand*{\definitionautorefname}{Definiton}
\newcommand*{\subfigureautorefname}{\figureautorefname}


%---- listings --------------------------------------------
% add \lineref command for listing enviroment
\newcommand*{\lineref}[2][]{%
\ifthenelse{\equal{#1}{}}{%
\hyperref[{#2}]{line~\ref*{#2}}%
}{%
\hyperref[{#2}]{\begin{NoHyper}line~\ref{#1}--\ref{#2}\end{NoHyper}}%
}\xspace%
}


%----- Beamer settings --------------------------------------
\IfClass{beamer}{
  \usetheme{rwth}
  \setbeamercovered{transparent}                                  % Versteckter Text ist transparent
  \setbeamertemplate{caption}{\insertcaption}             % Only text as caption
  \setbeamerfont{caption}{size={\footnotesize}}           % Caption size
}

% tweak footnote line
\renewcommand{\footnoterule}{%
  \kern -3pt
  \hrule width 1.5in
  \kern 2.6pt
}

% superscript footnote number
% \deffootnote{0em}{1.6em}{\textsuperscript{\thefootnotemark}}

\usepackage{scrpage2}


\setcounter{tocdepth}{2} % nur 2 Gliederungsebenen im Inhaltsverzeichnis

% ------ Fussnoten
\pagestyle{scrheadings} 

\refoot{\thispublishers ~- \textbf{\thistitle}}
\lofoot{Stand \thisdate ~- Version \thisversion}

\renewcommand*{\chapterpagestyle}{scrheadings}
% ------



% \loadglsentries[\acronymtype]{preamble/Acronyms}
% \makeglossaries                               % Make the glossary
\makeindex                                      % Make the index

\begin{document}

\begin{titlepage}

\begin{huge}
\bfseries\center
\thistitle

\end{huge}

\vspace{1em}

\begin{large}
\bfseries\noindent
%\thisauthor
\begin{tabular}{@{}ll@{}}
  Markus Beutel & RWTH Aachen University \\
  Sevket Gökay & RWTH Aachen University \\
  Peter von Grumbkow & HaCon Ing.-Ges. mbH \\
  Dirk Hillbrecht & Cantamen GmbH \\
  Karl-Heinz Krempels & RWTH Aachen University \\
  Christian Samsel & RWTH Aachen University \\
  Christoph Terwelp & RWTH Aachen University \\
  Heike Twele & HaCon Ing.-Ges. mbH \\
  Gerhard Wagner & Cantamen GmbH \\
\end{tabular}


\end{large}

\vspace{2em}

Datum: \thisdate ~ Dokumentversion: \thisversion

\vspace{2em}

%\begin{abstract}
\thisabstract
%\end{abstract}
\vspace{2em}
\vfill

\includegraphics[width=0.25\textwidth]{econnect-germany.pdf} \hspace{1.5em}
\includegraphics[width=0.35\textwidth]{mb-logo.png}
\hfill \includegraphics[width=0.20\textwidth]{BMWi_4C_M-Gef2.pdf}


\end{titlepage}


\cleardoublepage


\tableofcontents

% \listoffigures



\cleardoublepage

% \linenumbers

\chapter{Overview}
\label{cha:Zusammenfassung}

Aim of this interface specification is the information exchange of information systems for rentals vehicles with travel information systems.
The reason for the information exchange is the requirement of creating intermodal travel chains combining rental services and public transport.


The specification in-hand contains:
\begin{itemize}
\item A Role Model,
\item A Service hierachy of different qualities of information exchange,
\item Interaction sequences to explain the order of messages between the partners based on the different service levels,
\item The data model,
\item Recommendations for the use of specific technologies for the data exchange and parsing,
\item Tables of allowed values for enumerations.
\end{itemize}

\chapter{Rollenmodell}
\label{sec:Rollenmodell}

Das Rollenmodell beschreibt die vorkommenden Rollen des Integrationsszenarios eines Fahrzeugverleihsystems (hier CarSharing CS) in ein Reiseinformationssystem.

\section{Reiseinformationssystem}

\subsection{Beschreibung}
Das Reiseinformationssystem (RIS) stellt ein Informationssystem für Reiseauskünfte dar und umfasst die Zusammenführung der Reiseangebote, die Konstruktion von Reiseketten, die Berechnung des Gesamtpreises für eine Reisekette, die Reservierung der Mobilitätsangebote für die einzelnen Elemente einer Reisekette, die Aufbereitung der Daten für die Darstellung in Benutzerschnittstellen und abschließend die Darstellung der Informationen.

\section{Anwendungsfälle}
\begin{itemize}
\item Kunde fragt mit Suchbedingungen und Präferenzen Auskünfte über Mobilitätsangebote beim RIS an. Suchbedingungen und Präferenzen sind beispielsweise Start- und Zielort, Abfahrt- und Ankunftszeitpunkt, zu berücksichtigende Verkehrsmittel, Anzahl der Umstiege, Preisspanne, etc. Die Ergebnisse werden den Kunden über die Benutzerschnittstelle des RIS in Form von Reiseketten dargestellt. 
% \item Die Buchung einer Reisekette führt zur Erzeugung einer zeitlich und räumlich eingeschränkten Autorisierung  zur Nutzung der Mobilitätsangebote der gewählten Reisekette (Ticket) sowie der abschließenden Abrechnung der Dienstleistung.
\end{itemize}

\section{Fahrzeugverleihsystem}

\subsection{Beschreibung}
Das Fahrzeugverleihsystem (FVS) stellt ein Informationssystem zur Verwaltung und Buchung von Leihfahrzeugen und Kunden dar. Fahrzeuge können sowohl unterschiedlichen Typs, als auch Stationsgebunden oder Stationsungebunden sein.

\subsection{Anwendungsfälle}
\begin{itemize}
\item Ein Kunde bucht (leiht) über das FVS ein Fahrzeug zu vertraglich geregelten Preisen, Zeiten und Stationen aus und nutzt dieses.
\item Ein Kunde fragt über das FVS die Vefügbarkeit eines Fahrzeugs an. 
\end{itemize}

\section{Benutzer RIS}

\subsection{Beschreibung}
Kunde RIS – stellt eine juristische Person dar, die befugt ist eine Reise zu buchen und unter Nutzung der gewählten Verkehrsmittel, anzutreten. 

\subsection{Anwendungsfälle}
\begin{itemize}
\item Benutzer stellt Auskunftsanfrage an das RIS.
\item Benutzer bucht Reise über das RIS.
\end{itemize}

\section{Benutzer FVS}

\subsection{Beschreibung}
Kunde FVS  - stellt eine juristische Person dar, die befugt ist ein Fahrzeug zu leihen und zu nutzen.

\subsection{Anwendungsfälle}
\begin{itemize}
\item Benutzer stellt Auskunftsanfrage an das FVS.
% \item Benutzer reserviert ein Fahrzeug über das FVS.
\item Benutzer bucht ein Fahrzeug über das FVS.
\end{itemize}

\section{Betreiber RIS}

\subsection{Beschreibung}
Betreiber RIS – stellt das RIS als Informationssystem für einen Mobilitätsanbieter als Dienstleistung bereit.

% \subsection{Anwendungsfälle}

\section{Betreiber FVS}
\subsection{Beschreibung}
Betreiber FVS - stellt ein Informationssystem für FVS für einen Fahrzeugverleiher als Dienstleistung bereit.

% \subsection{Anwendungsfälle}

\section{FVS-Station}

\subsection{Beschreibung}
FVS-Station - stellt die Fahrzeuge eines Mobilitätsanbieters an einer Station bereit.

% \subsection{Anwendungsfälle}

\section{Fahrzeug}

\subsection{Beschreibung}
Fahrzeug – wird vom Mobilitätsanbieter über das FVS an einer Station als Mobilitätsangebot für den Kunden angeboten.

% \subsection{Anwendungsfälle}


\chapter{Hierarchical Model}
\label{cha:Hierachiemodell}
The hierarchical model describes different qualities of information coupling, based on service groups and serves as a recommendation concerning possible stages of implementation.
To realize a coupling between VRS and TIS, at least service 1 (static data) of both interacting parties has to be supported. Dependencies between different services are depicted in \cref{fig:depend}.

\begin{figure}[h]
  \centering
  \resizebox {\columnwidth} {!} {
  \begin{tikzpicture}
  % Dienste
  \tikzstyle{ann} = [draw=none,fill=none,right]
  \node[rectangleTX](s1) {Service 1\\Static data};
  \node[rectangleTX, below left of=s1](s3) {Service 3\\Availability Subscription};
  \node[rectangleTX, left= 1cm of s3](s2) {Service 2\\Availability Query};
  \node[rectangleTX, below right of=s1](s4) {Service 4\\Booking};
  \node[rectangleTX, right= 1cm of s4](s6) {Service 6\\Price Information};
  \node[rectangleTX, right= 1.5cm of s6](s9) {Service 9\\User Management};


  \node[rectangleTX, below = 2cm of s4](s8) {Service 8\\Booking status\\(vehicle unlock)};
  \node[rectangleTX, left= 0.5cm of s8](s7) {Service 7\\Consumption data};
  \node[rectangleTX, left= 0.5cm of s7](s5) {Service 5\\Booking Subscription};
  \node[rectangleTX, right= 0.5cm of s8](s10) {Service 10\\Vehicle Settings};
  \node[rectangleTX, right= 0.5cm of s10](s11) {Service 11\\Navigation and Assistance};
  %   arrows
  \draw (s2) edge[->,thick=1cm] (s1);
  \draw (s3) edge[->,thick=1cm] (s1);
  \draw (s4) edge[->,thick=1cm] (s1);
  \draw (s5) edge[->,thick=1cm] (s4);
  \draw (s6) edge[->,thick=1cm] (s1);
  \draw (s7) edge[->,thick=1cm] (s4);
  \draw (s8) edge[->,thick=1cm] (s4);
  \draw (s9) edge[->,thick=1cm] (s1);
  \draw (s10) edge[->,thick=1cm] (s4);
  \draw (s11) edge[->,thick=1cm] (s4);
  \end{tikzpicture}
  }
  \caption{Service Dependencies \label{fig:depend}}
\end{figure}


\section{Base Service A -- Session Handling}
\label{sec:Hierachiemodell:BasisdienstA}
Service A enables the authentication of end-customers towards the VRS.

\subsection*{Functions}
\begin{itemize}
\item Open / Close Session
\end{itemize}

\subsection*{Dependencies}
none

\section{Base Service B -- Subscription Handling}
\label{sec:Hierachiemodell:BasisdienstB}
Service B contains a function to check the status of a subscription connection (heartbeat)

\subsection*{Functions}
\begin{itemize}
\item Heartbeat
\end{itemize}

\subsection*{Dependencies}
none

\section{Base Service C -- Tokens}
\label{sec:Hierachiemodell:BasisdienstC}
Service C contains a function to create authentication tokens for users, which can be saved / transferred instead of plaintext passwords.

\subsection*{Functions}
\begin{itemize}
\item Creation of Tokens
\end{itemize}

\subsection*{Dependencies}
none


\section{Service 1 -- Static Data}
\label{sec:Hierachiemodell:Dienst1}
Service 1 serves the information exchange across vehicle rental companies and static data of booking targets. These include provider-, position- and vehicle-data.

For example, service 1 can be used, to display messages solely for existing locations of an VRS-provider in a TIS. 
\subsection*{Functions}
\begin{itemize}
\item Call of booking targets and provider information.
\end{itemize}

\subsection*{Dependencies}
none

\section{Service 2 -- Availability Query}
\label{sec:Hierachiemodell:Dienst2}
Service 2 serves for synchronous calls of availability information. The actual availability times of booking targets are called from the VRS during the travel inquiry through the TIS.

The actual availability times of booking targets are called during the travel inquiry by the TIS at the VRS.
\subsection*{Functions}
\begin{itemize}
\item Calls of availabilities of booking targets
\item Calls of location capacities (Service 2a)
\end{itemize}

\subsection*{Dependencies}
\begin{itemize}
\item Service 1
\end{itemize}

\section{Service 3 -- Availability Subscription}
\label{sec:Hierachiemodell:Dienst3}
Service 3 serves the asynchronous exchange of availability information.

To accelerate the travel inquiry, the TIS can subscribe availability timescales of booking targets to avoid a query during the travel inquiry. After subscribing for an amount of booking targets, the VRS informs continuously about changes in availability timescales. 

\subsection*{Functions}
\begin{itemize}
\item Availability subscription
\item Availability information (push)
\item Subscription of location capacities (Service 3a)
\end{itemize}

\subsection*{Dependencies}
\begin{itemize}
\item Service 1
\item Base Service B
\end{itemize}


\section{Service 4 -- Booking}
\label{sec:Hierachiemodell:Dienst4}
Service 4 serves booking, booking changes and canceling of vehicles via the TIS on behalf of the customer of a VRS.

The booking of a vehicle requires that a customer authenticates himself towards the VRS. Therefore, authentication information is forwarded from the TIS to the VRS. The secret (Password, key, etc.) of the customer is not allowed to be saved due to security reasons. The TIS receives an authentication-token in case of a successful authentication of a user. It can either be saved on the TIS or the user's device. Using the authentication-token, the TIS can proceed queries concerning bookings and canceling. 
To allow booking changes, a booking can be replaced by another booking via a booking change request. In case of a booking errors, it has to be ensured that the initial booking retains it's validity. 

\subsection*{Functions}
\begin{itemize}
\item Authentication of users towards the VRS 
\item Booking Query
\item Query concerning booking changes/ canceling
\item Subscription of external Bookings (Service 4a)
\end{itemize}

\subsection*{Dependencies}
\begin{itemize}
\item Service 1
\end{itemize}

\section{Service 5 -- Booking Subscription}
\label{sec:Hierachiemodell:Dienst5}
Service 5 serves the subscription of booking changes.

The TIS is able to subscribe conducted bookings at the VRS to inform users in case of changes, e.g., damaged vehicles.
\subsection*{Functions}
\begin{itemize}
\item Booking subscriptions 
\item Booking alert (Push)
\end{itemize}

\subsection*{Dependencies}
\begin{itemize}
\item Service 4
\item Base Service B
\end{itemize}

\section{Service 6 -- Price Information}
\label{sec:Hierachiemodell:Dienst6}
Service 6 serves the price information of rental services.

Through the transmission of starting- and target-location as well as departure- and arrival-time of the travel, the TIS is able to query price information at the VRS to inform the user. The VRS responds with an overall price and eventually individual price parts.

\subsection*{Functions}
\begin{itemize}
\item Query of Prices
\end{itemize}

\subsection*{Dependencies}
\begin{itemize}
\item Service 1
\end{itemize}


\section{Service 7 -- Subscription Consumption Data/ Billing}
Service 7 serves to exchange consumption data, e.g, usage duration and distance. Through the transmission of consumption data, the TIS is able to create bills for end-customers. Therefore, the TIS subscribes consumption data of an executed booking at the VRS and receives the corresponding consumption data automatically. 

\subsection*{Functions}
\begin{itemize}
\item Subscription consumption data 
\item Exchange consumption data (Push)
\end{itemize}

\subsection*{Dependencies}
\begin{itemize}
\item Service 4
%\ item Basisdienst A
\item Base service B
\end{itemize}

\section{Service 8 -- Booking State Change (Unlocking / Pausing / Finishing)}
Service 8 serves to unlock booking of respective vehicles (probably via key case) by the TIS. This allows to unlock as well as to return of the vehicle during the usage phase of the TIS, e.g. via a mobile application. Furthermore, a booking can be paused, which locks the vehicle physically without finishing the booking. 
The state always refers to a preceding booking and an authentication is required.  

\subsection*{Functions}
\begin{itemize}
\item Unlock booking
\item Pause booking 
\item Finish booking 
\end{itemize}

\subsection*{Dependencies}
\begin{itemize}
\item Service 4
\item Base service A
\end{itemize}

\section{Service 9 -- User Management}
Service 9 serves to create and lock of user accounts of the VRS through the TIS. Therefore, a user just has to register himself at the TIS operator. This registration is valid even for usage of the VRS.

\subsection*{Functions}
\begin{itemize}
\item Create user account
\item Lock user account
\end{itemize}

\subsection*{Dependencies}
\begin{itemize}
\item Service 1
\item Base service A
\end{itemize}


\section{Service 10 -- Vehicle or Booking Settings}
Service 10 serves the management of vehicle settings. The TIS is able to transfer vehicle settings (e.g., air condition temperature, preferred radio stations) to the VRS and subscribe to settings of existing bookings.

\subsection*{Functions}
\begin{itemize}
\item Set vehicle settings 
\item Subscribe to vehicle settings
\end{itemize}

\subsection*{Dependencies}
\begin{itemize}
\item Service 4
\item Base service A
\end{itemize}

\section{Service 11 -- Remote Configuration of Navigation System and Journey Progress Monitoring}
Service 11 serves the integration of the vehicle navigation system into an intermodal journey. The TIS is able to transfer the target location to the vehicle after the booking, so that the traveler does not need to enter the location by his self. Additionally, the TIS is able to monitor the progress of a journey, in order to react in case of delays.

\subsection*{Functions}
\begin{itemize}
\item Set navigation target of vehicle
\item Subscribe to route progress
\end{itemize}

\subsection*{Dependencies}
\begin{itemize}
\item Service 4
\item Base service A
\end{itemize}

\chapter{Interaktionsprotokolle}
\label{sec:Interaktionsprotokolle}

\section{Ebene 0 -- Stationsinformationen}

\subsection{Handshake}
\interactiongraphics{Level0Handshake}
Um die Version und die Ebene des Protokolls mit der das RIS und das FVS miteinander kommunizieren festzulegen, wird direkt nach dem Verbindungsaufbau ein Handshake durchgeführt. Dabei teilt das FVS dem Reiseinformationssystem, die unterstützten Protokollversionen und die maximale unterstützte Protokollebene mit. Das RIS wählt daraus die gewünschte Dienstebene aus und teilt sie dem FVS mit.

\subsection{Abfrage der Betreiber}
\interactiongraphics{Level0Betreiber}
Da über eine Verbindung Informationen über mehrere FV Betreiber ausgetauscht werden können, ist es notwendig, dass die beteiligten Betreiber von einander unterschieden werden können und zu diesen jeweils grundlegende Informationen bereitgestellt werden können. Hierfür wird eine Betreiberliste durch das FVS bereitgestellt, deren Aktualität durch das FVS anhand einer Version-ID festgestellt werden kann. Die Versions-ID darf vom RIS nur gespeichert und für spätere Anfrage verwendet werden. Eine anderweitige Verwendung der Versions-ID durch das RIS ist nicht vorgesehen. 

Als Parameter der Abfrage muss die Versions-ID verwendet werden, falls bereits eine Betreiberliste vorliegt.
Wenn bei einer Anfrage durch das RIS eine Versions-ID angegeben wurde, muss das FVS prüfen, ob sich die Versions-ID der aktuellen Betreiberliste des FVS von der gegebenen Versions-ID unterscheidet. Falls sich die Versions-IDs unterscheiden oder keine Versions-ID als Parameter angegeben wurde, antwortet das FVS mit der aktuellen Betreiberliste und ihrer Versions-ID. Falls die Versions-IDs gleich sind, informiert das FVS das RIS darüber, dass sich die Betreiberliste nicht verändert hat.

\subsection{Abfrage der Stationen}
\interactiongraphics{Level0Stationen}
Die Stationen eines FVS beschreiben Standorte an denen Fahrzeuge ausgeliehen und zurückgegeben werden können. Eine Liste von Stationen bezieht sich immer auf genau einen Betreiber, welcher in einer Anfrage anhand seiner Betreiber-ID identifiziert wird. Es ist möglich die Auswahl der Stationen durch die Angabe eines PLZ-Bereichs einzuschränken. Analog zur Abfrage der Betreiber wird eine Versions-ID verwendet.
Die Antwort des FVS muss genau alle Stationen enthalten, die den Anfragekriterien (Betreiber-ID, PLZ-Bereich) entsprechen.

\section{Ebene 1 -- Verfügbarkeitsauskunft}

\subsection{Abfrage der Zielgebiete}
\interactiongraphics{Level1Zielgebiete}
Zielgebiete dienen der Auskunft eines Free-Floating-Fahrzeugverleihs. Ein Zielgebiet ist ein Bereich, in denen ein Verleihvorgang eines Fahrzeugs beendet werden kann.
Eine Liste von Zielgebieten bezieht sich immer auf genau einen Betreiber, der bei der Anfrage durch seine Betreiber-ID identifiziert wird. Analog zur Abfrage der Betreiber wird eine Versions-ID verwendet. Die Antwort des FVS muss genau alle Zielgebiete des angegebenen Betreibers enthalten.

\subsection{Abfrage der Buchungsziele}
\interactiongraphics{Level1Buchungsziel}
Um eine Auskunft über die zur Verfügung stehenden Fahrzeuge zu realisieren, werden die Fahrzeuge als Buchungsziele abgebildet. Es handelt sich hierbei nicht zwingend um eine eins-zu-eins Beziehung zwischen Fahrzeug und Buchungsziel. Die Anfrage der Buchungsziele muss durch die Angabe einer Betreiber-ID auf einen Betreiber eingeschränkt erfolgen. Optional können analog zum Abruf der Stationsinformationen ein PLZ-Bereich oder eine Stations-ID (im Fall eines Betreibers von Free-Floating-Fahrzeugeverleih, ein Polygon bestehend aus GPS-Positionen) zur Filterung der Buchungsziele angegeben werden. Das FVS antwortet zunächst mit einer vollständigen Liste aller den Filterkriterien entsprechenden Buchungsziele und einer für das FVS eindeutigen Abo-ID. Die Abo-ID wird für die Zuordnung nachfolgender Aktualisierungsnachrichten verwendet. Buchungsziele werden durch ihre Buchungsziel-ID zeitunabhängig eindeutig identifiziert.
Buchungsziele sind im Gegensatz zu Stations- und Betreiberinformationen dynamische Informationen, die sich häufig ändern. Daher werden diese kontinuierlich aktualisiert.
Bei der Änderung eines Buchungsziels im FVS wird dieses erneut an das RIS übertragen, wobei die Buchungsziel-ID übereinstimmen muss.
Zur Löschung eines Buchungsziels wird eine Löschmitteilung mit der Buchungsziel-ID übertragen.
Buchungsziele, die zur Ergebnisliste hinzukommen, werden mit einer neuen eindeutigen Buchungsziel-ID übertragen. Wenn für eine Anfrage keine weiteren Aktualisierungen benötigt werden, kann diese unter Angaben der Abo-ID beendet werden.
 
\section{Ebene 2 -- Tarifauskunft}

\subsection{Abfrage der Tarifliste}
\interactiongraphics{Level2Tarife}
Die Tarife eines FVS beschreiben vereinfachte Vertragskonditionen zu denen Fahrzeuge aus Buchungszielen angemietet werden können.
Eine Liste von Tarifen bezieht sich immer auf genau einen Betreiber, welcher in einer Anfrage anhand seiner Betreiber-ID identifiziert wird. Analog zur Abfrage der Betreiber wird eine Versions-ID verwendet.
Die Antwort des FVS muss genau die Tarife des gewählten Betreibers enthalten.
Werden nur die Dienstebenen 0 bis 2 der Schnittstelle verwendet, wird nur die Liste der öffentlich zugänglichen Tarife vom FVS dem RIS bereitgestellt.

\section{Ebene 3 -- Buchung}

\subsection{Authentifizierung eines Benutzers}
\interactiongraphics{Level3Authentifizierung}
Buchungen durch das RIS im Namen eines Nutzers erfordern dessen Authentifizierung beim FVS.
Hierfür wird eine Betreiber-ID zur Auswahl des Betreibers, ein Benutzeridentifikationsmerkmal (z.B. Benutzername, Email, ...) und das Passwort des Benutzers benötigt. Das FVS prüft diese Informationen auf Korrektheit und erzeugt bei Erfolg ein Authentifizierungs-Token. Dieses wird bei allen weiteren benutzerbezogenen Operation zur Authentifizierung des Benutzers verwendet. Die sicherere Speicherung des Authentifizierungs-Tokens liegt in der Verantwortung des RIS. Es kann auf dem Endgerät des Benutzers oder im Benutzerprofil im RIS gespeichert werden.

\subsection{Abfrage einer Preisauskunft}
\interactiongraphics{Level3Preisauskunft}
Eine Preisauskunft durch das RIS und die preisoptimierte Konstruktion von intermodalen Reiseketten erfordern eine Auskunft durch das FVS über den zu erwartenden Preis der Buchung eines Fahrzeugs.
Dazu werden bei der Anfrage neben der Betreiber-ID, die Buchungsziel-ID, ein Zeitfenster für die geplante Buchung, die Tarif-ID des anzuwendenden Tarifs, die erwartete Wegstrecke und das Identifikationsmerkmal des Benutzers angegeben. Das FVS antwortet mit einer Preisauskunft, in welcher der zu erwartende Preis aufgeteilt ist, in die Komponenten Grundbetrag, Zeitkosten, Distanzkosten und sonstige Kosten angegeben.
 
\subsection{Anfrage zur genauen Verfügbarkeit}
\interactiongraphics{Level3Verfuegbarkeit}
Buchungsziele müssen nicht eins-zu-eins jedes Fahrzeug des FVS Betreibers abbilden, sondern können auch nur eine maximale Verfügbarkeit darstellen. Daher können die Buchungsziele nur genutzt werden um sicher festzustellen, wann kein Fahrzeug mit bestimmten Parametern zur Verfügung steht. Um festzustellen, ob ein Fahrzeug zur Verfügung steht, wird eine Anfrage zur Verfügbarkeit eines Fahrzeugs an das FVS gestellt. Bei der Anfrage werden die Betreiber-ID, Buchungsziel-ID, ein Zeitfenster für welche die Verfügbarkeit geprüft werden soll und das Identifikationsmerkmal des Benutzers für den die Auskunft genutzt werden soll angegeben. Im Fall von Stationsflexiblem oder Free-Floating FV wird eine Zielstation bzw. ein Zielbereich durch ihre entsprechend ID angegeben. Wenn für die gegebenen Parameter ein Fahrzeug zur Verfügung steht antwortet das FVS mit einer Zeitspanne für die ein Fahrzeug zur Verfügung und einer Liste der anwendbaren Tarife. Die Zeitspanne muss größer oder gleich der bei der Anfrage angegebenen Zeitspanne sein. Falls kein Fahrzeug zur Verfügung steht, informiert das FVS das RIS in der Antwort.

\subsection{Auftrag zur Buchung}
\interactiongraphics{Level3Buchung}
Ein Buchungsauftrag enthält zusätzlich zu allen Informationen die auch bei einer Verfügbarkeitsanfrage übertragen werden das Benutzer-Token und die Tarif-ID des gewählten Tarifs. Das FVS prüft die Durchführbarkeit der Buchung und antwortet nach durchgeführter Buchung mit den Buchungsinformationen oder andernfalls mit einer Fehlermeldung, die Aufschluss über den Grund gibt, warum die Buchung nicht durchgeführt werden konnte.

\subsection{Anfrage für Aktualisierungen zu einer Buchung}
\interactiongraphics{Level3Buchungsaktualisierungen}
Anpassungen einer Reisekette durch das RIS bei Änderungen einer Buchung erfordern die Übertragung von Buchungsänderungen vom FVS zum RIS. Änderungen einer Buchung können durch das RIS beim FVS abonniert werden. Dazu sendet das RIS eine Anfrage mit der Betreiber-ID und Buchungs-ID an das FVS. Diese wird durch das FVS mit einer für das FVS eindeutigen Abo-ID und dem aktuellen Stand der Buchung bestätigt. Nachfolgende Änderungen der Buchung werden unter Verwendung der Abo-ID durch das FVS an das RIS gesendet. Wenn das Ende der Buchung erreicht wurde, die Buchung storniert wurde oder das RIS eine Anfrage zur Beendigung des Abonnements mit der Abo-ID, werden folgende Änderungen der Buchung nicht mehr an das RIS übertragen.

\subsection{Auftrag zur Stornierung einer Buchung}
\interactiongraphics{Level3Stornierung}
Der Auftrag zur Stornierung einer Buchung enthält die Betreiber-ID, die Buchungs-ID der zu stornierenden Buchung und zur Authentifizierung das Benutzer-Token. Das FVS bestätigt dem RIS die Durchführung der Stornierung.

\section{Ebene 4 -- Umbuchung}

\subsection{Auftrag zur Umbuchung}
\interactiongraphics{Level4Umbuchung} 
Eine Umbuchung erlaubt in einem Schritt eine Buchung zu verändern oder durch eine neue Buchung zu ersetzen. Dabei muss durch das FVS sichergestellt werden, dass bei einem Fehler die ursprüngliche Buchung bestehen bleibt.
\chapter{Datenmodell}
\label{sec:Datenmodell}

\section{Dienst 1}
\label{subsec:Datenmodell:Dienst1}

\subsection*{Buchungsziel}

\input{xml/generated/BookingTargetType-docu}
% \TODO{XML Beispiel}

\subsection*{Verleihstation}

\input{xml/generated/PlaceType-docu}
% \TODO{XML Beispiel}

\subsection*{Provider}

\input{xml/generated/ProviderType-docu}
% \TODO{XML Beispiel}

\subsection*{Fahrzeugeigenschaft}

\input{xml/generated/AttributeType-docu}
% \TODO{XML Beispiel}


\section{Dienst 2}
\label{subsec:Datenmodell:Dienst1}

\subsection*{Buchungszieleigenschaften}

\input{xml/generated/BookingTargetPropertiesType-docu}
% \TODO{XML Beispiel}

\subsection*{Zeitrahmen}

\input{xml/generated/TimePeriodType-docu}
% \TODO{XML Beispiel}
\chapter{Technische Realisierung}
\label{sec:TechnischeRealisierung}


\section{Nachrichtenkodierung}
The messages are transferred as XML-files between both systems. A precise type-definition is set by the XML-scheme, corresponding to the interface.

In case the overhead, resulting of the introduction of XML turns problematic, there is the possibility to use the Efficient XML Interchange (EXI) protocol. The use of EXI would reduce the size of messages significantly, without loosing the advantages of XML. 

\section{Kommunikationskanal}
As the interface allows asynchronous subscription besides the usual request/response scheme, usage of the WebSocket protocol instead of plain HTTP is recommended.
WebSockets allow persistent connections between both systems and a bidirectional message exchange. Because the interface also provides an asynchronous subscription model, besides the request and response scheme and a low response rate is desirable, the WebSocket-protocol is recommended for the communication channel. WebSockets allow to establish a connection between two systems and exchange messages bidirectional. The VRS represents the server and the TIS the client (in HTTP context). In principal, any amount of communication channels can be opened. Updates of subscribed objects are delivered via the same connection as the one they are subscribed at. In case of an interrupt of the connection, the subscription ends.

\section{Authentifizierung}
IXSI is a B2B interface and therefore does not include an internal authentication mechanism. Instead the communication partners are advised to use existing authentication mechanisms such as SSL certificates, a virtual private network (VPN) or HTTP authentication.

\subsection*{Endkunde}
Because via IXSI, even user-controlled requests are transferred from system to system, it is necessary that a user can authenticate himself towards the VRS.
This usually happens via the tripel provider reference/user reference/password. To avoid typing the password in plaintext, alternatively to the password, a token can be used. This token can be generated by the VRS with the help of the function  TokenGeneration (vgl. \cref{subsec:Nachrichten:DienstC:Tokengenerierung}).
Thus, the user initially has to type his password once, afterwards he is able to authenticate himself via the generated token. For example, this can be saved on the user's device. 

In case of various consecutive requests which are invoked by the same user, a authentication should not be executed in every one. Therefore, with a user's first request, a session will be opened (explicitly via OpenSession or implicitly). For following requests of the same user, the ID of the opened session instead of an authentication can be transferred. These session have a time limited validity. In case this validity expires, the user needs to be identified again and a new session needs to be opened. In addition, a session can be closed explicitly via CloseSession.

\section{Verbindungssicherheit}
To ensure confidentiality of the transferred data, encryption is required. Therefore, the SSL/TLS-protocol is suitable. This should be used in case that the connection is not already secured through other measures (e.g., by using VPN).

\chapter{Nachrichten}
\label{cha:Nachrichten}
Die zwischen den Interaktionspartnern ausgetauschten Nachrichten basieren auf fünf Nachrichtentypen der Art Request / Response und Push. Alle Nachrichten vom RIS an das FVS stellen Anfragen dar und sind daher vom Typ Request. Der Typ Response wird für direkte Antworten des FVS auf Anfragen des RIS verwendet. Um eine Zuordnung der Response-Nachrichten zu den zugehörigen Request-Nachrichten zu erlauben, wird jede Request-Nachricht mit einer eindeutigen Transaction-ID markiert, die in der zugehörigen Response-Nachricht wieder mitgegeben werden muss. Eine Response-Nachricht kann anstelle ihres normalen Inhalts auch einen Fehler enthalten. Der letzte Nachrichtentyp Push findet in den Fällen Verwendung, wenn das RIS fortlaufende Aktualisierungen (Abonnement) FVS angefordert hat. 
\TODO{Beispiel Nachrichten einfügen}

\section{Basisnachrichten}
\label{sec:Nachrichten:Basisnachrichten}

\subsection*{Basisklasse}
\input{xml/generated/IxsiMessageType}
\input{xml/generated/IxsiMessageType-schema}

\subsection*{Basisantwort}
\input{xml/generated/AbstractBaseResponseType}
\input{xml/generated/AbstractBaseResponseType-schema}
\medskip

\textit{Eine abstrakte Basisklasse für einen Request ist nicht vorhanden, da nicht erforderlich.}

\subsection*{Transaktionsdaten}
\input{xml/generated/TransactionType}
\input{xml/generated/TransactionType-schema}

\subsection*{Query Request}
\input{xml/generated/QueryRequestType}
\input{xml/generated/QueryRequestType-schema}

\subsection*{Query Response}
\input{xml/generated/QueryResponseType}
\input{xml/generated/QueryResponseType-schema}

\subsection*{Subscription Request}
\input{xml/generated/SubscriptionRequestType}
\input{xml/generated/SubscriptionRequestType-schema}

\subsection*{Subscription Response}
\input{xml/generated/SubscriptionResponseType}
\input{xml/generated/SubscriptionResponseType-schema}

\section{Basisgruppen}
\label{sec:Nachrichten:Basisgruppen}

\subsection*{Benutzer-initiert}
\input{xml/generated/UserTriggeredRequestGroup}
\input{xml/generated/UserTriggeredRequestGroup-schema}
\input{xml/generated/UserTriggeredResponseGroup}
\input{xml/generated/UserTriggeredResponseGroup-schema}

\subsection*{Herkunft/Zielangabe}
\input{xml/generated/OriginDestGroup}
\input{xml/generated/OriginDestGroup-schema}

\subsection*{Gruppen rein zur Organisation}
Anmerkung: Im XML Schema sind weitere Gruppen definiert die jedoch nur der Organisation dienen und keinen Einfluss auf das Nachrichtenformat haben und deshalb hier nicht aufgelistet sind.
\begin{samepage}
\begin{flushleft}
\rowcolors{1}{}{gray!10}
\begin{tabularx}{\linewidth}{l>{\raggedright\arraybackslash}X} 
\toprule
Name &  Kommentar \label{tab:miscGroups}\\
\midrule 
\texttt{StaticDataRequestGroup} & \\
\texttt{StaticDataResponseGroup}& \\
\texttt{SubscriptionAdministrationRequestGroup} & \\
\texttt{SubscriptionAdministrationResponseGroup}& \\
\texttt{SubscriptionRequestGroup} & \\
\texttt{SubscriptionResponseGroup}& \\
\texttt{RequestMessageGroup}& \\
\texttt{ResponseMessageGroup}& \\
\texttt{PushMessageGroup}& \\
\bottomrule 
\end{tabularx}\end{flushleft}\end{samepage}

\section{Dienst 1}
\label{sec:Nachrichten:Dienst1}

\subsection*{Abfrage Buchungsziele}
\label{subsec:Nachrichten:Dienst1:BookingTargets}
\input{xml/generated/BookingTargetsInfoRequestType}
\input{xml/generated/BookingTargetsInfoRequestType-schema}
\input{xml/generated/BookingTargetsInfoResponseType}
\input{xml/generated/BookingTargetsInfoResponseType-schema}

\subsection*{Abfrage Änderungen Buchungszielen}
\label{subsec:Nachrichten:Dienst1:ChangeProviders}
\input{xml/generated/ChangedProvidersRequestType}
\input{xml/generated/ChangedProvidersRequestType-schema}
\input{xml/generated/ChangedProvidersResponseType}
\input{xml/generated/ChangedProvidersResponseType-schema}

\section{Dienst 2}
\label{sec:Nachrichten:Dienst2}

\subsection*{Abfrage Verfügbarkeiten}
\label{subsec:Nachrichten:Dienst2:Availability}
\input{xml/generated/AvailabilityRequestType}
\input{xml/generated/AvailabilityRequestType-schema}
\input{xml/generated/AvailabilityResponseType}
\input{xml/generated/AvailabilityResponseType-schema}

\subsection*{Abfrage Standortkapazitäten (Dienst 2a)}
\label{subsec:Nachrichten:Dienst2:PlaceAvailability}
\input{xml/generated/PlaceAvailabilityRequestType}
\input{xml/generated/PlaceAvailabilityRequestType-schema}
\input{xml/generated/PlaceAvailabilityResponseType}
\input{xml/generated/PlaceAvailabilityResponseType-schema}

\section{Dienst 3}
\label{sec:Nachrichten:Dienst3}

\subsection*{Verfügbarkeitsabonnement}
\label{subsec:Nachrichten:Dienst3:AvailabilitySubscription}
\input{xml/generated/AvailabilitySubscriptionRequestType}
\input{xml/generated/AvailabilitySubscriptionRequestType-schema}
\input{xml/generated/AvailabilitySubscriptionResponseType}
\input{xml/generated/AvailabilitySubscriptionResponseType-schema}
\input{xml/generated/AvailabilityPushMessageType}
\input{xml/generated/AvailabilityPushMessageType-schema}

\subsection*{Vollständige Verfügbarkeitsabonnementinformation}
\label{subsec:Nachrichten:Dienst3:CompleteAvailability}
\input{xml/generated/CompleteAvailabilityRequestType}
\input{xml/generated/CompleteAvailabilityRequestType-schema}
\input{xml/generated/CompleteAvailabilityResponseType}
\input{xml/generated/CompleteAvailabilityResponseType-schema}

\subsection*{Standortkapazitätsabonnement (Dienst 3a)}
\label{subsec:Nachrichten:Dienst3:PlaceAvailabilitySubscription}
\input{xml/generated/PlaceAvailabilitySubscriptionRequestType}
\input{xml/generated/PlaceAvailabilitySubscriptionRequestType-schema}
\input{xml/generated/PlaceAvailabilitySubscriptionResponseType}
\input{xml/generated/PlaceAvailabilitySubscriptionResponseType-schema}
\input{xml/generated/PlaceAvailabilityPushMessageType}
\input{xml/generated/PlaceAvailabilityPushMessageType-schema}

\subsection*{Vollständige Standortkapazitätsinformation (Dienst 3a)}
\label{subsec:Nachrichten:Dienst3:CompletePlaceAvailability}
\input{xml/generated/CompletePlaceAvailabilityRequestType}
\input{xml/generated/CompletePlaceAvailabilityRequestType-schema}
\input{xml/generated/CompletePlaceAvailabilityResponseType}
\input{xml/generated/CompletePlaceAvailabilityResponseType-schema}

\section{Dienst 4}
\label{sec:Nachrichten:Dienst4}

\subsection*{Buchung}
\label{subsec:Nachrichten:Dienst4:Booking}
\input{xml/generated/BookingRequestType}
\input{xml/generated/BookingRequestType-schema}
\input{xml/generated/BookingResponseType}
\input{xml/generated/BookingResponseType-schema}

\subsection*{Buchungsänderung}
\label{subsec:Nachrichten:Dienst4:ChangeBooking}
\input{xml/generated/ChangeBookingRequestType}
\input{xml/generated/ChangeBookingRequestType-schema}
\input{xml/generated/ChangeBookingResponseType}
\input{xml/generated/ChangeBookingResponseType-schema}

\section{Dienst 5}
\label{sec:Nachrichten:Dienst5}

\subsection*{Buchungsabonnement}
\label{subsec:Nachrichten:Dienst5:BookingAlertSubscription}
\input{xml/generated/BookingAlertSubscriptionRequestType}
\input{xml/generated/BookingAlertSubscriptionRequestType-schema}
\input{xml/generated/BookingAlertSubscriptionResponseType}
\input{xml/generated/BookingAlertSubscriptionResponseType-schema}
\input{xml/generated/BookingAlertPushMessageType}
\input{xml/generated/BookingAlertPushMessageType-schema}

\subsection*{Vollständige Buchungsabonnementinformation}
\label{subsec:Nachrichten:Dienst5:CompleteBookingAlert}
\input{xml/generated/CompleteBookingAlertRequestType}
\input{xml/generated/CompleteBookingAlertRequestType-schema}
\input{xml/generated/CompleteBookingAlertResponseType}
\input{xml/generated/CompleteBookingAlertResponseType-schema}

\section{Dienst 6}
\label{sec:Nachrichten:Dienst6}

\subsection*{Preisinformation}
\label{subsec:Nachrichten:Dienst6:PriceInformation}
\input{xml/generated/PriceInformationRequestType}
\input{xml/generated/PriceInformationRequestType-schema}
\input{xml/generated/PriceInformationResponseType}
\input{xml/generated/PriceInformationResponseType-schema}

\section{Dienst 7}
\label{sec:Nachrichten:Dienst7}

\subsection*{Nutzungsdatenabonnement}
\label{subsec:Nachrichten:Dienst7:ConsumptionSubscription}
\input{xml/generated/ConsumptionSubscriptionRequestType}
\input{xml/generated/ConsumptionSubscriptionRequestType-schema}
\input{xml/generated/ConsumptionSubscriptionResponseType}
\input{xml/generated/ConsumptionSubscriptionResponseType-schema}
\input{xml/generated/ConsumptionSubscriptionStatusRequestType}
% \input{xml/generated/ConsumptionSubscriptionStatusRequestType-schema}
\input{xml/generated/ConsumptionSubscriptionStatusResponseType}
% \input{xml/generated/ConsumptionSubscriptionStatusResponseType-schema}
\input{xml/generated/ConsumptionPushMessageType}
\input{xml/generated/ConsumptionPushMessageType-schema}

\subsection*{Vollständige Nutzungsdaten}
\label{subsec:Nachrichten:Dienst5:CompleteConsumption}
\input{xml/generated/CompleteConsumptionRequestType}
\input{xml/generated/CompleteConsumptionRequestType-schema}
\input{xml/generated/CompleteConsumptionResponseType}
\input{xml/generated/CompleteConsumptionResponseType-schema}

\section{Dienst A}
\label{sec:Nachrichten:DienstA}

\subsection*{Session öffnen}
\label{subsec:Nachrichten:DienstA:OpenSession}
\input{xml/generated/OpenSessionRequestType}
\input{xml/generated/OpenSessionRequestType-schema}
\input{xml/generated/OpenSessionResponseType}
\input{xml/generated/OpenSessionResponseType-schema}

\subsection*{Session schließen}
\label{subsec:Nachrichten:DienstA:CloseSession}
\input{xml/generated/CloseSessionRequestType}
\input{xml/generated/CloseSessionRequestType-schema}
\input{xml/generated/CloseSessionResponseType}
\input{xml/generated/CloseSessionResponseType-schema}

\section{Dienst B}
\label{sec:Nachrichten:DienstB}

\subsection*{Heartbeat}
\label{subsec:Nachrichten:DienstB:HeartBeat}
\input{xml/generated/HeartBeatRequestType}
\input{xml/generated/HeartBeatRequestType-schema}
\input{xml/generated/HeartBeatResponseType}
\input{xml/generated/HeartBeatResponseType-schema}

\section{Dienst C}
\label{sec:Nachrichten:DienstC}

\subsection*{Tokengenerierung}
\label{subsec:Nachrichten:DienstC:Tokengenerierung}
\input{xml/generated/TokenGenerationRequestType}
\input{xml/generated/TokenGenerationRequestType-schema}
\input{xml/generated/TokenGenerationResponseType}
\input{xml/generated/TokenGenerationResponseType-schema}

\chapter{Code Tabellen}
\label{cha:CodeTabellen}
Die syntaktische Definition für IXSI enthält keine konkreten Werte (Enumerations) für beispielsweise Fahrzeugtypen oder Fehlercodes. Stattdessen werden diese Werte in den nachfolgenden Codetabellen festgelegt. Nur die hier festgelegten Werte dürfen innerhalb von IXSI verwendet werden.

\section{Fahrzeugklassen}
\label{sec:CodeTabellen:ClassType}
Für die Auswahl \emph{ClassType} können folgende Werte verwendet werden:

\begin{flushleft}
\rowcolors{1}{}{gray!10}
\begin{tabularx}{\linewidth}{l>{\raggedright\arraybackslash}X}
\toprule
Wert & Bedeutung\\
\midrule
\verb|bike| & Fahrrad\\
\verb|motorcycle| & Motorrad\\
\verb|micro| & Kleinstwagen (z.\,B. Smart4two)\\
\verb|mini| & Kleinwagen (z.\,B. Opel Corsa)\\
\verb|small| & Kompaktwagen (z.\,B. VW Golf)\\
\verb|medium| & Mittelklassewagen (z.\,B. Audi A4)\\
\verb|large| & Oberklassewagen (z.\,B. BMW 7er)\\
\verb|van| & Van (z.\,B. VW T5 Multivan)\\
\verb|transporter| & Transporter (z.\,B. Ford Transit)\\
\bottomrule
\end{tabularx}
\end{flushleft}

\section{Antriebsklassen}
\label{sec:CodeTabellen:EngineType}
Für die Auswahl \emph{EngineType} können folgende Werte verwendet werden:

\begin{flushleft}
\rowcolors{1}{}{gray!10}
\begin{tabularx}{\linewidth}{l>{\raggedright\arraybackslash}X}
\toprule
Wert & Bedeutung\\
\midrule
\verb|none| & Kein Kraftantrieb (Muskelkraft)\\
\verb|diesel| & Dieselmotor\\
\verb|gasoline| & Ottomotor\\
\verb|electric| & Elektromotor\\
\verb|liquidgas| & Flüssiggas (LPG)\\
\verb|naturalgas| & Erdgas (CNG)\\
\verb|hydrogen| & Wasserstoffantrieb\\
\verb|hybrid| & Hybridantrieb mit Elektro- und Verbrennungsmotor\\
\bottomrule
\end{tabularx}
\end{flushleft}

Hinweis: Ein Pedelec kann als Kombination Fahrzeugklasse \verb|bike| und Antriebsklasse \verb|electric| dargestellt werden.

\section{Fahrzeugeigenschaften}
\label{sec:CodeTabellen:AttributeClass}
Attribute und Eigenschaften eines Buchungsziels oder eines Standorts können klassifiziert werden, um sie automatisch interpretieren zu können. Dazu werden Attributsklassen verwendet. Attributsklassen werden in IXSI im Datentyp \emph{AttributeClassType} gespeichert. Folgende Werte sind dabei erlaubt:

\begin{flushleft}
\rowcolors{1}{}{gray!10}
\begin{tabularx}{\linewidth}{l>{\raggedright\arraybackslash}X}
\toprule
Wert & Bedeutung\\
\midrule
\verb|trailer_hitch| & Anhängerkupplung\\
\verb|automatic| & Automatikgetriebe\\
\verb|convertible| & Cabriolet\\
\verb|air_condition| & Klimaanlage\\
\verb|navigation| & Navigationssystem\\
\verb|cruise_control| & Tempomat\\
\verb|winter_tyres| & Winter- bzw. Ganzjahresreifen\\
\verb|child_seat_0| & Babyschale\\
\verb|child_seat_1| & Kindersitz (9-18kg)\\
\verb|child_seat_4| & Kindersitz (15-36kg)\\
\verb|utility| & Kombi\\
\verb|doors_4| & 4/5-Türer\\
\verb|seats_9| & Mindestens 9 Sitze\\
\verb|seats_7| & Mindestens 7 Sitze\\
\verb|seats_5| & Mindestens 5 Sitze\\
\verb|seats_4| & Mindestens 4 Sitze\\
\bottomrule
\end{tabularx}
\end{flushleft}

\section{Verbrauchsdaten}
\label{sec:CodeTabellen:ConsumptionClassType}
Klassen von Verbrauchsdaten zur Abrechnung von Leistungen, wie z.\,B. Entfernung, Dauer oder Zusatzleistungen wie ein Kindersitz. Wird in  \emph{ConsumptionClassType} gespeichert. Folgende Werte sind dabei erlaubt:

\begin{flushleft}
\rowcolors{1}{}{gray!10}
\begin{tabularx}{\linewidth}{l>{\raggedright\arraybackslash}X}
\toprule
Wert & Bedeutung\\
\midrule
\verb|distance| & Entfernung\\
\verb|duration| & Dauer\\
\bottomrule
\end{tabularx}
\end{flushleft}

\section{Nutzerzustände}
\label{sec:CodeTabellen:UserStateType}

Zustände, die ein Benutzer im System haben kann, wird in  \emph{UserStateType} gespeichert. Folgende Werte sind dabei erlaubt:

\begin{flushleft}
\rowcolors{1}{}{gray!10}
\begin{tabularx}{\linewidth}{l>{\raggedright\arraybackslash}X}
\toprule
Wert & Bedeutung\\
\midrule
\verb|operative| & Benutzer ist aktiviert und darf Fahrzeuge buchen.\\
\verb|nonoperative| & Benutzer ist gesperrt und darf keine Fahrzeuge buchen (kann entsperrt werden).\\
\verb|deleted| & Benutzer ist (endgültig) gelöscht.\\
\bottomrule
\end{tabularx}
\end{flushleft}
\section{Fehlercodes}
\label{sec:CodeTabellen:ErrorCode}
Das FVS kann alternativ zu den normalen Response-, Update- und Handshake-Nachrichten auch im Fall eines Fehlers Fehlernachrichten an das RIS schicken. Fehlercodes werden im Datentyp \emph{ErrorCodeType} gespeichert.
\begin{flushleft}
\rowcolors{1}{}{gray!10}
\begin{tabularx}{\linewidth}{l>{\raggedright\arraybackslash}X}
\toprule
Wert & Bedeutung\\
\midrule
\verb|auth_provider_unknown| & Authentifizierung: Unbekannte Provider-ID\\
\verb|auth_invalid_password| & Authentifizierung: User-Passwort-Kombination ungültig\\
\verb|auth_invalid_token| & Authentifizierung: User-Token-Kombination ungültig\\
\verb|auth_session_invalid| & Authentifizierung: Session ist ungültig/abgelaufen\\
\verb|auth_anon_not_allowed| & Authentifizierung: Anonymer User nicht erlaubt\\
\verb|auth_not_authorized| & Autorisation: Nutzer ist zu dieser Anfrage nicht berechtigt\\
\verb|sys_backend_failed| & System: Hintergrundsystem antwortet nicht\\
\verb|sys_unknown_failure| & System: unbekannter Fehler\\
\verb|sys_not_implemented| & System: Request nicht implementiert\\
\verb|sys_request_not_plausible| & System: Request ist nicht plausibel. Dieser Wert sollte stets verwendet werden, wenn inhaltlich Fehler im Request enthalten sind.\\
\verb|booking_target_unknown| & Buchungsziel unbekannt\\
\verb|price_info_not_available| & Preisinformationen nicht verfügbar\\
\verb|booking_too_short| & Buchungsdauer zu kurz\\
\verb|booking_too_long| & Buchungsdauer zu lang\\
\verb|booking_target_not_available| & Buchungsziel im gegebenen Zeitraum nicht buchbar\\
\verb|booking_change_not_possible| & Buchungsänderung kann nicht durchgeführt werden\\
\verb|booking_id_unknown| & Unbekannte Buchungs-ID. Dieser Wert sollte auch verwendet werden, wenn die Buchungs-ID einem anderen User zugeordnet ist.\\
\verb|booking_state_change_failure| & Änderung des Buchungsstatus fehlgeschlagen.\\
\verb|language_not_supported| & Angefragte Sprache nicht vollständig unterstützt, andere Sprache geliefert.\\
\bottomrule
\end{tabularx}
\end{flushleft}


\cleardoublepage

\appendix


% \setglossarysection{section}
% \printglossary[type=\acronymtype,style=long]

\renewcommand{\lstlistlistingname}{Quellcodeverzeichnis}
\addcontentsline{toc}{chapter}{\lstlistlistingname}
\lstlistoflistings

\renewcommand{\indexname}{Stichwortverzeichnis}
\printindex

% \bibliographystyle{plain}
%{\small\bibliography{literature/Literature}}


\vfill
\end{document}

